\documentclass[10pt, aspectratio=169]{beamer}
\usepackage{../beamer_theme_408}

\title{408考研零基础 --- 计算机组成原理}
\subtitle{2-17 总线系统}
\author{}
\date{}

\begin{document}

\frame{\titlepage}

\begin{frame}{本节目标}
    \begin{enumerate}
        \item 理解总线的定义与基本特性
        \item 掌握总线的分类方式
        \item 了解三种集中式总线仲裁方式
        \item 区分同步总线与异步总线的时序
    \end{enumerate}
\end{frame}

\begin{frame}{总线的定义与特性}
    \textbf{总线}：一组能为多个部件\textbf{分时共享}的公共信息传送线路
    \vspace{0.3em}
    \begin{block}{两个关键特征}
        \begin{itemize}
            \item \textbf{共享}：多个部件连接到同一组总线上
            \item \textbf{分时}：同一时刻只能有一个部件发送数据
        \end{itemize}
    \end{block}
    \vspace{0.3em}
    \textbf{总线的特性}：
    \begin{itemize}
        \item \textbf{机械特性}：尺寸、引脚数、引脚排列
        \item \textbf{电气特性}：信号电平、传输方向
        \item \textbf{功能特性}：各引脚的功能定义
        \item \textbf{时间特性}：信号的时序关系
    \end{itemize}
\end{frame}

\begin{frame}{总线的分类}
    \begin{enumerate}
        \item \textbf{按数据传输方式}
            \begin{itemize}
                \item 并行总线（多根数据线同时传多位）
                \item 串行总线（单根数据线逐位传输）
            \end{itemize}
        \item \textbf{按使用范围}
            \begin{itemize}
                \item 片内总线（CPU芯片内部）
                \item 系统总线（CPU、主存、I/O之间）
                \item 通信总线（计算机与外部设备之间）
            \end{itemize}
        \item \textbf{按系统总线内容}
            \begin{itemize}
                \item 数据总线（DB）：传输数据信息，双向
                \item 地址总线（AB）：传输地址信息，单向
                \item 控制总线（CB）：传输控制信号，双向
            \end{itemize}
    \end{enumerate}
\end{frame}

\begin{frame}{系统总线的三大组成部分}
    \begin{center}
        \begin{tabular}{|c|c|c|c|}
            \hline
            \textbf{总线类型} & \textbf{传输内容} & \textbf{方向} & \textbf{位宽} \\
            \hline
            数据总线 DB & 数据 / 指令 & 双向 & 8/16/32/64位 \\
            地址总线 AB & 内存地址 / I/O端口 & 单向（CPU输出） & 决定寻址空间 \\
            控制总线 CB & 读/写/中断等控制信号 & 双向 & 多根控制线 \\
            \hline
        \end{tabular}
    \end{center}
    \vspace{0.5em}
    \begin{block}{重要关系}
        \begin{itemize}
            \item 地址总线宽度 $\rightarrow$ 最大寻址空间（$2^n$，$n$为地址线根数）
            \item 数据总线宽度 $\rightarrow$ 一次传送的数据位数
        \end{itemize}
    \end{block}
\end{frame}

\begin{frame}{总线仲裁 --- 为什么需要仲裁？}
    \begin{itemize}
        \item 总线是共享资源
        \item 同一时刻只能有一个主设备占用总线
        \item 多个主设备同时请求时，需\textbf{仲裁}决定谁先使用
    \end{itemize}
    \vspace{0.5em}
    \textbf{集中式仲裁}的三种方式：
    \vspace{0.3em}
    \begin{enumerate}
        \item \textbf{链式查询}：按优先顺序串联
        \item \textbf{计数器定时查询}：由计数器轮询
        \item \textbf{独立请求方式}：每个设备独立请求线
    \end{enumerate}
\end{frame}

\begin{frame}{链式查询方式}
    \begin{center}
        \includegraphics[width=0.7\textwidth]{chain_query.png}
    \end{center}
    \vspace{0.3em}
    \begin{itemize}
        \item BR（总线请求）线：所有设备共用
        \item BG（总线允许）线：逐级传递
        \item 离仲裁器最近的设备优先级最高
    \end{itemize}
    \vspace{0.3em}
    \begin{tabular}{|c|c|}
        \hline
        \textbf{优点} & \textbf{缺点} \\
        \hline
        电路简单，线数少 & 优先级固定 \\
        易扩展 & 单点故障（一级断，全部瘫痪） \\
        \hline
    \end{tabular}
\end{frame}

\begin{frame}{计数器定时查询与独立请求}
    \textbf{计数器定时查询}
    \begin{itemize}
        \item 仲裁器用一个计数器轮询各设备
        \item 从计数值开始查，查到请求即停
        \item 优先级可编程（改变计数初值）
        \item 比链式查询灵活，但线数较多
    \end{itemize}
    \vspace{0.3em}
    \textbf{独立请求方式}
    \begin{itemize}
        \item 每个设备有独立的请求线和允许线
        \item 仲裁器内部并行判决
        \item 速度最快，优先级最灵活
        \item 但线数最多（$2n$根）
    \end{itemize}
\end{frame}

\begin{frame}{总线的操作时序}
    \textbf{同步总线}
    \begin{itemize}
        \item 由统一的时钟信号控制数据传输
        \item 所有设备按固定时钟节拍工作
        \item 优点：时序简单、实现容易
        \item 缺点：时钟频率受限于最慢设备
    \end{itemize}
    \vspace{0.3em}
    \textbf{异步总线}
    \begin{itemize}
        \item 没有统一时钟，采用\textbf{握手协议}
        \item 请求 $\rightarrow$ 应答 $\rightarrow$ 完成
        \item 优点：能适配不同速度的设备
        \item 缺点：时序复杂，效率低于同步
    \end{itemize}
\end{frame}

\begin{frame}{常用总线标准}
    \begin{center}
        \begin{tabular}{|c|c|c|c|}
            \hline
            \textbf{总线标准} & \textbf{类型} & \textbf{数据宽度} & \textbf{用途} \\
            \hline
            PCI & 同步、并行 & 32/64位 & 系统内部扩展 \\
            PCIe & 异步、串行 & 1/4/8/16通道 & 高速外围设备 \\
            USB & 异步、串行 & 1位 & 外部设备连接 \\
            SATA & 异步、串行 & 1位 & 硬盘接口 \\
            \hline
        \end{tabular}
    \end{center}
    \vspace{0.5em}
    \begin{itemize}
        \item 现代总线趋势：串行代替并行（更高频率、更少干扰）
        \item PCIe采用差分信号传输，速度快、抗干扰强
    \end{itemize}
\end{frame}

\begin{frame}{本节总结}
    \begin{enumerate}
        \item 总线是计算机各部件间分时共享的公共通信线路
        \item 总线分为片内总线、系统总线、通信总线
        \item 系统总线含数据、地址、控制三组总线
        \item 集中式仲裁：链式查询、计数器查询、独立请求
        \item 同步总线用统一时钟，异步总线用握手协议
    \end{enumerate}
    \vspace{1em}
    \begin{center}
        \textcolor{blue}{\textbf{下节预告：2-18 I/O系统概述 --- CPU如何与外设通信}}
    \end{center}
\end{frame}

\end{document}
